\section{Metodología}
La arquitectura del proyecto se diseñó modularmente mediante un orquestador \texttt{main.py} que ejecuta secuencialmente las cuatro etapas.

\subsection{Etapa 1: Preprocesamiento}
Los conjuntos de datos \textit{Slice.txt} y \textit{Vh.txt} fueron convertidos a un formato estandarizado \texttt{.csv} (\textit{Comma Separated Values}) para facilitar la carga vectorial mediante la biblioteca Pandas.

\subsection{Etapa 2: Normalización}
La normalización se aplicó a las matrices de características de ambos datasets, generando tres versiones transformadas por archivo: \textit{Min-Max}, \textit{Z-Score} y \textit{Robust}. Posteriormente, se seleccionó empíricamente la versión Z-Score como línea base para la reducción de dimensionalidad por su eficacia en evitar sesgos de escala.

\subsection{Etapa 3: Análisis y Reducción}
A partir del dataset estandarizado con Z-Score, se ejecutaron tres procesos paralelos de reducción:
\begin{enumerate}
    \item \textbf{Selección de Características:} Se calculó la matriz de covarianza usando la correlación de Pearson. Recorriendo la matriz triangular superior, se identificaron y eliminaron aquellas características con una correlación $|r| > 0.98$.
    \item \textbf{Proyección PCA:} Se aplicó PCA para reducir el conjunto a 10 componentes principales, registrando el porcentaje de varianza conservada.
    \item \textbf{Transformación CCA:} Se aplicó CCA para maximizar la correlación entre las características continuas y las etiquetas discretas.
\end{enumerate}

\subsection{Etapa 4: División de Datos}
Para prevenir la filtración de datos (\textit{Data Leakage}) durante el futuro entrenamiento, los datos reducidos se partieron utilizando un muestreo estratificado (\texttt{stratify=y}) para preservar la distribución de clases original. La división fue: $70\%$ Entrenamiento (\textit{Train}), $15\%$ Validación (\textit{Validation}) y $15\%$ Prueba (\textit{Test}).